\documentclass[11pt]{article}
\usepackage{amsmath}
\usepackage{amssymb}
\usepackage{amsthm}
\usepackage{enumitem}
\usepackage[margin=1in]{geometry}
\usepackage{multicol}
\usepackage{soul}
\usepackage{tikz}

% Load hyperref to automatically generate the PDF outline from structural tags
\usepackage[pdftex, bookmarks=true, bookmarksnumbered=true]{hyperref}

% Optional: Use the bookmark package for faster rendering and advanced control
\usepackage{bookmark} 

\newenvironment{case}[2]{\par\textbf{Case #1: #2.}}{}
\title{Homework \#4\\{\large Math 104: Intro to Analysis \\ Summer 2026} \\ {\normalsize UC Berkeley}}
\date{July 24, 2026}
\author{Donald Aingworth IV}

\begin{document}
\renewcommand\qedsymbol{OE$\Delta$}
\newcommand{\abs}[1]{\left| #1 \right|}
\newcommand{\andtxt}{\text{ and }}
\newcommand{\Ais}{\ensuremath{\overset{A}{=}}}
\newcommand{\Cis}{\ensuremath{\overset{C}{=}}}
\newcommand{\closure}[1]{\ensuremath{\overline{#1}}}
\newcommand{\encircle}[1]{%
  \tikz[baseline=(char.base)]{%
    \node[shape=circle, draw, inner sep=1pt] (char) {#1};%
  }%
}
\newcommand{\mathif}{\text{if }}
\newcommand{\seq}[2][n]{\ensuremath{\left\{ #2 _{#1} \right\}}}
\newcommand{\zz}[1]{$\mathbb{Z}/ #1 \mathbb{Z}$}

\maketitle

\tableofcontents

\pagebreak
\setcounter{section}{-1}
\section{Various Theorems}
    \subsection{Monotone Decrease I}
        $a_n = \frac{1}{n(\log n)(\log \log n)^p}$ for $n \in \mathbb{N}$ is monotonically decreasing if $p > 0$.
        \begin{proof}
            Let $p > 0$.
            We know $n > 0$ since $n \in \mathbb{N}$.
            This also means $n \geq 1$.
            Logarithms are monotone, as mentioned in Rudin (p.63).
            % Rudin also mentions that $\frac{1}{n\log(n)}$ converges to $0$.
            Therefore, $\log n$ and by extension $\log \log n$ and $(\log \log n)^p$ are monotonically increasing.
            Therefore, $n (\log n) (\log \log n)^p$ is monotonically increasing.
            On the other hand, this also means that $\frac{1}{n (\log n) (\log \log n)^p}$ is monotonically decreasing.
        \end{proof}

    \subsection{Monotone Decrease II}
        $b_n = \frac{1}{n (\log 2) (\log (n) + \log (\log 2))^p}$ for $n \in \mathbb{N}$ is monotonically decreasing if $p > 0$.
        \begin{proof}
            This uses a lot of what was used in Monotone Decrease I.
            $n$ is monotone increasing.
            $\log(n)$ is monotone increasing and $\log(\log(2))$ is constant, so $\log(n) + \log(\log(2))$ and by extension $(\log(n) + \log(\log(2)))^p$ if $p > 0$ are monotone increasing.
            $\log(2)$ is constant, so $n (\log 2) (\log (n) + \log (\log 2))^p$ is monotone increasing if $p > 0$.
            Therefore, $\frac{1}{n (\log 2) (\log (n) + \log (\log 2))^p}$ is monotone decreasing if $p > 0$.
        \end{proof}

    \subsection{Inverese Triangle Inequality}
        For $a, b, c \in (X,d)$, $\abs{d(p,q) - d(q,r)} \leq d(p,r)$.
        \begin{proof}
            From the triangle inequality, we have two equations.
            \begin{gather}
                d(p,q) \leq d(p,r) + d(q,r)\\
                d(q,r) \leq d(p,r) + d(p,q)
            \end{gather}

            We can subtract $d(q,r)$ and $d(p,q)$ from them respectively.
            \begin{gather}
                d(p,q) - d(q,r) \leq d(p,r)\\
                d(q,r) - d(p,q) \leq d(p,r)
            \end{gather}

            Reverse the latter.
            \begin{gather}
                d(p,q) - d(q,r) \geq -d(p,r)
            \end{gather}

            By the definition of the absolute value, we can say that $\abs{d(p,q) - d(q,r)} \leq d(p,r)$.
        \end{proof}

    \subsection{Finite Times Natural is Finite}
        A finite number times any natural number is finite.
        \begin{proof}[Proof by Induction]
            Let $x \in \mathbb{R}$ be finite, so $x \neq \pm \infty$.

            Base case: $n = 1$.
            If $n = 1$, then $n \cdot x = 1 \cdot x = x$.
            Since $x$ is finite, this means the base case holds.

            Inductive step: suppose that $nx$ is finite.
            Take $(n + 1)x$.
            The distributive property says that $(n + 1)x = nx + x$.
            Since both $x$ and $nx$ are finite, $nx + x$ is finite.
            Ergo, $(n + 1)x$ is finite.
        \end{proof}

\pagebreak
\section{Problem 1}
    For which values of $p \in \mathbb{R}$ are the following series convergent or divergent? Prove your answers rigorously, using only tests we explicitly covered in class. \textit{Hint:} for both of these series you should use something with the name ``Cauchy" in it.
    \begin{enumerate}[label=(\alph*)]
        \item 
        $\displaystyle{\sum_{n = 3}^{\infty} \frac{1}{n(\log n)(\log \log n)^p}}$

        \item 
        $\displaystyle{\sum_{n = 1}^{\infty}} \frac{(-1)^{n - 1}}{n^p}$.
    \end{enumerate}

    \subsection{Solution (a)}
    % First we prove that $\sum_{n=1}^{\infty} \frac{1}{n(\log n)(\log \log n)^p}$ converges if and only if $\sum_{n=1}^{\infty} \frac{1}{(\log 2) (n \log (2) + \log (\log 2))^p}$ converges.
    \begin{proof}
        Since $2 < e$, $\log(2) < \log(e) = 1$.
        Furthermore, through the use of a calculator, $\log 2 = 0.6931$ and $\log \log 2 = -0.3665$.

        Let $p \leq 0$.
        This means $a_n = \frac{(\log \log n)^q}{n (\log n)}$ for some $q \geq 0$ such that $q = -p$.
        $\log n$ is monotonically increasing, so by extension so is $\log \log n$.
        Therefore, since $\log n$ is monotonically increasing and unbounded, so is $\log \log n$ and so is $(\log \log n)^q$.
        Therefore, for some $N_0 \in \mathbb{N}$, for all $n \geq N_0$, $\log \log n \geq 1$ and $(\log \log n)^q \geq 1$.
        Therefore if $n \geq N_0$, $\frac{(\log \log n)^q}{n \log n} \geq \frac{1}{n \log n}$.
        Recall that if $a_n \geq d_n \geq 0$ for $n \geq N_0$ and if $\sum d_n$ diverges, then $\sum a_n$ diverges (Rudin, p.60).
        In Rudin (p.62), it is proven that $\sum_{n=2}^{\infty} \frac{1}{n\log n}$ diverges.
        Therefore, $\sum_{n = 3}^{\infty} \frac{(\log \log n)^q}{n \log n}$ diverges.

        Going forward, we will therefore only operate on cases where $p > 0$.
        We established that if $p > 0$, there exists $N_0 \in \mathbb{N}$ such that for all $n \geq N_0$, $\frac{1}{n(\log n)(\log \log n)^p} > 0$\footnote{Might need further proof, but see prior paragraph.}.
        Furthermore, it is established earlier in this document that $\frac{1}{n(\log n)(\log \log n)^p}$ is monotonically decreasing.

        We will use the Cauchy Condensation Test, which says that for monotonically decreasing, positive $a_n$, $\sum_{n=1}^{\infty} a_n$ converges if and only if $\sum_{n=1}^{\infty} 2^n a_{2^n}$ converges.
        Consider it in this case.
        \begin{gather}
            a_n =   \frac{1}{n(\log n)(\log \log n)^p}\\
            b_n =   2^n a_{2^n} 
                =   \frac{2^n}{2^n(\log 2^n)(\log \log 2^n)^p}
                =   \frac{1}{n (\log 2) (\log (n \log 2))^p}
                =   \frac{1}{n (\log 2) (\log (n) + \log (\log 2))^p}
        \end{gather}

        Now, we can use how $\frac{1}{n \log 2 (\log (n) + \log (\log 2))^p}$ is positive and monotonically decreasing to use the Cauchy Condensation Test again.
        \begin{gather}
            b_n =   \frac{1}{n (\log 2) (\log (n) + \log (\log 2))^p}\\
            c_n =   2^n b_{2^n}
                =   \frac{2^n}{2^n (\log 2) (\log (2^n) + \log (\log 2))^p}
                =   \frac{1}{(\log 2) (n \log (2) + \log (\log 2))^p}
        \end{gather}

        Here, we have two cases: either $p > 1$ or $p \leq 1$.

        First consider the case where $p \leq 1$ where $n \in \mathbb{N}$ and $n \geq 3$.
        Now, we know that $\log(2) < 1$, so $n\log(2) < n$.
        Therefore, since $\log\log 2 < 0$, $n\log(2) + \log \log (2) < n$ and by extension $(n\log(2) + \log \log (2))^p < n^p$.
        Therefore, $\frac{1}{(n\log(2) + \log \log (2))^p} > \frac{1}{n^p}$.
        If we sum this for $n \geq 3$, we find that $\sum_{n=3}^{\infty} \frac{1}{(n\log(2) + \log \log (2))^p} > \sum_{n=3}^{\infty} \frac{1}{n^p}$.
        Since $\sum_{n=1}^{\infty} \frac{1}{n^p}$ diverges to $\infty$ for $p \leq 1$, we can equally say that $\sum_{n=3}^{\infty} \frac{1}{n}$ diverges to $\infty$ under the same conditions.
        Recall again that if $a_n \geq d_n \geq 0$ for $n \geq N_0$ and if $\sum d_n$ diverges, then $\sum a_n$ diverges (Rudin, p.60).
        Since $\sum_{n=3}^{\infty} \frac{1}{n^p}$ diverges to $\infty$ for $p \leq 1$, $\sum_{n=3}^{\infty} \frac{1}{(n\log(2) + \log \log (2))^p}$ also diverges to $\infty$ for $p \leq 1$.
        Therefore, by the Cauchy Condensation Test, if $0 < p \leq 1$, $\sum_{n = 3}^{\infty} \frac{1}{n(\log n)(\log \log n)^p}$ diverges.

        Now we consider the case where $p > 1$ for $n \in \mathbb{N}$ and $n \geq 3$.
        I will contend right now that for all $n \in \mathbb{N}$ where $n \geq 3$, $n\log(2) + \log \log (2) > \frac{n}{2}$ and will prove it inductively.
        The base case is that $n = 3$.
        With the help of our calculator, we can find that $3\log(2) + \log \log (2) = 1.7129 > \frac{3}{2}$, so our base case holds.
        For the inductive step, we assume that $n\log(2) + \log \log (2) > \frac{n}{2}$.
        Our calculator tells us that $\log(2) = 0.693$, so $\log(2) > \frac{1}{2}$.
        Since $a > b$ and $c > d$ implies that $a + c > b + d$, $n\log(2) + \log \log (2) + \log(2) > \frac{n}{2} + \frac{1}{2}$.
        Therefore, $(n + 1)\log(2) + \log \log (2) > \frac{n + 1}{2}$, so the inductive step holds.
        If $p > 1$, then $\left( n\log(2) + \log \log (2) \right)^p > \left( \frac{n}{2} \right)^p$.
        Therefore, $\frac{1}{\left( n\log(2) + \log \log (2) \right)^p} < \frac{2^p}{n^p}$.

        We know that if $p > 1$, then $\sum_{n = 1}^{\infty} \frac{1}{n^p}$ is convergent and by extension $\sum_{n = 3}^{\infty} \frac{1}{n^p}$.
        Since a series is still a sequence of partial sums, $2^p$ is finite if $p \in \mathbb{R}$ and $p > 1$, and a convergent sequence times a finite constant converges, $2^p \sum_{n = 3}^{\infty} \frac{1}{n^p} = \sum_{n = 3}^{\infty} \frac{2^p}{n^p}$ still converges.
        If $\abs{a_n} \leq c_n$ for all $n \in \mathbb{N}$ and $n \geq N_0$ for some $N_0 \in \mathbb{Z}$ and $\sum c_n$ converges, then $a_n$ converges (Rudin, p.60).
        Therefore, since $\frac{1}{\left( n\log(2) + \log \log (2) \right)^p}$ is made up of only terms that are always positive, it is always positive, so $\abs{\frac{1}{\left( n\log(2) + \log \log (2) \right)^p}} = \frac{1}{\left( n\log(2) + \log \log (2) \right)^p}$.
        Therefore, $\abs{\frac{1}{\left( n\log(2) + \log \log (2) \right)^p}} \leq \frac{2^p}{n^p}$ for all $n \in \mathbb{N}$ and $n \geq 3$.
        Since $\sum_{n = 3}^{\infty} \frac{2^p}{n^p}$ converges, $\sum_{n = 3}^{\infty} \frac{1}{\left( n\log(2) + \log \log (2) \right)^p}$ also converges.
        Therefore, by the Cauchy Condensation Test, if $p > 1$, $\sum_{n = 3}^{\infty} \frac{1}{n(\log n)(\log \log n)^p}$ converges.

        Ergo, $\sum_{n = 3}^{\infty} \frac{1}{n(\log n)(\log \log n)^p}$ converges if $p > 1$ and diverges otherwise.
    \end{proof}

    \subsection{Solution (b)}

\pagebreak
\section{Problem 2}
    A series $\sum_{n = 1}^{\infty} a_n$ is said to be \textit{absolutely convergent} if $\sum_{n = 1}^{\infty} |a_n|$ is convergent. Prove that every absolutely convergent series is convergent. \textit{Hint:} use the Cauchy criterion.
    \begin{proof}
        Suppose that $\sum_{n=1}^{\infty} a_n$ is absolutely convergent.
        This means that $\sum_{n=1}^{\infty} \abs{a_n}$ is convergent and by extension Cauchy.
        By the Cauchy Criterion, this means that for all $\varepsilon > 0$, there exists some $N \in \mathbb{N}$ such that if $m \geq n \geq N$, $\abs{\sum_{k=n}^{m} \abs{a_k}} \leq \varepsilon$.

        Recall the triangle inequality for the usual metric: if $a, b, c \in \mathbb{R}$, $\abs{a - c} \leq \abs{a - b} + \abs{b - c}$.
        Since $\abs{a - c} = \abs{a - b + b - c}$, $\abs{a - b + b - c} \leq \abs{a - b} + \abs{b - c}$.
        Suppose that for some $a_k = a - b$ and $a_{k+1} = b - c$.
        Therefore, $\abs{a_k + a_{k + 1}} \leq \abs{a_k} + \abs{a_{k+1}}$.
        Since $m - n$ is finite, we can extend this process to every $k \in [n,m-1] \cap \mathbb{N}$ to conclude that $\abs{\sum_{k=n}^{m} a_k} \leq \sum_{k=n}^{m} \abs{a_k}$.
        By the absolute value, $\abs{\sum_{k=n}^{m} a_k} \geq 0$, meaning $\sum_{k=n}^{m} \abs{a_k} \geq 0$ and therefore $\sum_{k=n}^{m} \abs{a_k} = \abs{\sum_{k=n}^{m} \abs{a_k}}$.
        This means that $\abs{\sum_{k=n}^{m} a_k} \leq \abs{\sum_{k=n}^{m} \abs{a_k}}$.

        Transistively, we can say from our earlier convergence statement that for all $\varepsilon > 0$, there exists some $N \in \mathbb{N}$ such that if $m \geq n \geq N$, $\abs{\sum_{k=n}^{m} a_k} \leq \varepsilon$.
        Therefore, by the Cauchy Criterion, $\sum_{n = 1}^{\infty} a_n$ is convergent.
        Ergo, if a series is absolutely convergent, it is convergent.
    \end{proof}

\pagebreak
\section{Problem 3}
    Let $X, Y$ be sets and $f : X \to Y$ a function.
    \begin{enumerate}[label=(\alph*)]
        \item 
        Show that, for $E, F \subset Y$ we have $f^{-1}(E \cup F) = f^{-1}(E) \cup f^{-1}(F)$ and $f^{-1}(E \cap F) = f^{-1}(E) \cap f^{-1}(F)$.

        \item 
        Show that, for $E, F \subset X$, we have $f(E \cup F) = f(E) \cup f(F)$. Is it also always true that $f(E \cap F) = f(E) \cap f(F)$? Prove or give a counterexample.

        \item 
        Let $E \subset X$. Show that $E \subset f^{-1}(f(E))$. Show that equality holds for all $E$ if and only if $f$ is injective.

        \item 
        Let $E \subset Y$. Show that $f(f^{-1}(E)) \subset E$. Show that equality holds for all $E$ if and only if $f$ is surjective.
        
    \end{enumerate}

    \subsection{Solution (a)}
        First we prove that $f^{-1}(E \cup F) = f^{-1}(E) \cup f^{-1}(F)$.
        \begin{proof}
            Begin with $f^{-1}(E \cup F)$.
            If $x \in f^{-1}(E \cup F)$, then $f(x) \in E \cup F$.
            Therefore, $f(x) \in E$ or $f(x) \in F$.
            If $f(x) \in E$, then $x \in f^{-1}(E)$.
            If $f(x) \in F$, then $x \in f^{-1}(F)$.
            Therefore, since either $f(x) \in E$ or $f(x) \in F$, $x \in f^{-1}(E)$ or $x \in f^{-1}(F)$.
            Therefore, $x \in f^{-1}(E) \cup f^{-1}(F)$.
            This means $f^{-1}(E \cup F) \subset f^{-1}(E) \cup f^{-1}(F)$.

            For the reverse direction, begin with $x \in f^{-1}(E) \cup f^{-1}(F)$.
            This means that either $x \in f^{-1}(E)$ or $x \in f^{-1}(F)$.
            Without loss of generality, suppose that $x \in f^{-1}(E)$.
            This means that, $f(x) \in E$.
            Quite logically, this means that $f(x) \in E$ or $f(x) \in F$.
            Therefore, $f(x) \in E \cup F$.
            Therefore, $x \in f^{-1}(E \cup F)$.
            Concluding, $f^{-1}(E) \cup f^{-1}(F) \subset f^{-1}(E \cup F)$.

            Since $f^{-1}(E \cup F) \subset f^{-1}(E) \cup f^{-1}(F)$ and $f^{-1}(E) \cup f^{-1}(F) \subset f^{-1}(E \cup F)$, $f^{-1}(E \cup F) = f^{-1}(E) \cup f^{-1}(F)$.
        \end{proof}

        Next, we prove that $f^{-1}(E \cap F) = f^{-1}(E) \cap f^{-1}(F)$.
        \begin{proof}
            Begin with $f^{-1}(E \cap F)$.
            If $x \in f^{-1}(E \cap F)$, then $f(x) \in E \cap F$.
            Therefore, $f(x) \in E$ and $f(x) \in F$.
            Since $f(x) \in E$, $x \in f^{-1}(E)$.
            Since $f(x) \in F$, $x \in f^{-1}(F)$.
            Therefore, $x \in f^{-1}(E) \cap f^{-1}(E)$.
            Therefore, $f^{-1}(E \cap F) \subset f^{-1}(E) \cap f^{-1}(F)$.

            Suppose $x \in f^{-1}(E) \cap f^{-1}(F)$.
            Then, $x \in f^{-1}(E)$ and $x \in f^{-1}(F)$.
            Since $x \in f^{-1}(E)$, $f(x) \in E$.
            Since $x \in f^{-1}(F)$, $f(x) \in F$.
            Therefore, $f(x) \in E \cap F$.
            Therefore, $x \in f^{-1}(E \cap F)$.
            Therefore, $f^{-1}(E) \cap f^{-1}(F) \subset f^{-1}(E \cap F)$.

            Since $f^{-1}(E \cap F) \subset f^{-1}(E) \cap f^{-1}(F)$ and $f^{-1}(E) \cap f^{-1}(F) \subset f^{-1}(E \cap F)$, $f^{-1}(E \cap F) = f^{-1}(E) \cap f^{-1}(F)$.
        \end{proof}

    \subsection{Solution (b)}
        First, we prove that $f(E \cup F) = f(E) \cup f(F)$.
        \begin{proof}
            Suppose that $x \in X$ and $f(x) \in f(E \cup F)$.
            This means that $x \in E \cup F$.
            Therefore, either $x \in E$ or $x \in F$.
            Without loss of generality, suppose that $x \in E$.
            Therefore, $f(x) \in f(E)$.
            Using the same steps, we can show that $f(x) \in f(F)$.
            Therefore, $f(x) \in f(E) \cup f(F)$.
            This means that $f(E \cup F) \subset f(E) \cup f(F)$.

            Suppose that $x \in X$ and $f(x) \in f(E) \cup f(F)$.
            This means that either $f(x) \in f(E)$ or $f(x) \in f(F)$.
            Therefore either $x \in E$ or $x \in F$.
            This means that $x \in E \cup F$.
            Therefore, $f(x) \in f(E \cup F)$.
            This means that $f(E) \cup f(F) \subset f(E \cup F)$.

            Since and $f(E \cup F) \subset f(E) \cup f(F)$ and $f(E) \cup f(F) \subset f(E \cup F)$, $f(E \cup F) = f(E) \cup f(F)$.
        \end{proof}

        Now we disprove that $f(E \cap F) = f(E) \cap f(F)$.
        \begin{proof}[Proof by Counterexample]
            Take the case of $E = (0,\infty)$ and $F = (-\infty, 0)$ where $f: \mathbb{R} \to \mathbb{R}$ and $f(x) = \abs{x}$.
            Taking the intersection, $E \cap F = \varnothing$.
            Therefore, $f(E \cap F) = \varnothing$.

            Now, applying the function to $E$ and $F$, we find $f(E) = (0,\infty)$ and $f(F) = (0,\infty)$.
            Taking the intersection of the two, we find $f(E) \cap f(F) = (0,\infty)$.
            
            Since $(0,\infty) \neq \varnothing$, $f(E \cap F) \neq f(E) \cap f(F)$.
        \end{proof}

    \subsection{Solution (c)}
        First, we prove that $E \subset f^{-1}(f(E))$.
        \begin{proof}
            Let $x \in E$.
            Therefore, $f(x) \in f(E)$.
            Since $f(x) \in f(E)$ and $f^{-1}(f(E)) = \{x \in X \Big| f(x) \in f(E)\}$, $x \in f^{-1}(f(E))$.
            Therefore, $E \subset f^{-1}(f(E))$.
        \end{proof}

        Next we prove that if $f$ is injective, then $E = f^{-1}(f(E))$.
        \begin{proof}
            Injective means that if $a, b \in E$, $f(a) = f(b)$ implies that $a = b$.
            Suppose that $f$ is injective.
            Recall the definitions that $f(E) = \left\{f(x) \Big| x \in E\right\}$ and $f^{-1}(f(E)) = \left\{ x \Big| f(x) \in f(E) \right\}$.

            Take any element in $x \in f^{-1}(f(E))$.
            By the definition of $f^{-1}(f(E))$, this means that $f(x) \in f(E)$.
            By the definition of $f(E)$, there exists $y \in E$ such that $f(x) = f(y)$.
            Since $f$ is injective, this means that $x = y$.
            Therefore, since $y \in E$ and $x = y$, $x \in E$.
            This means that if $x \in f^{-1}(f(E))$, $x \in E$.
            Therefore, $f^{-1}(f(E)) \subset E$.

            We proved earlier that $E \subset f^{-1}(f(E))$.
            Therefore, $f^{-1}(f(E)) \subset E$ and $E \subset f^{-1}(f(E))$.
            Combining them, that means $E = f^{-1}(f(E))$.
        \end{proof}

        Third, we prove that if $E = f^{-1}(f(E))$, then $f$ is injective.
        \begin{proof}
            Recall the definitions that $f(E) = \left\{f(x) \Big| x \in E\right\}$ and $f^{-1}(f(E)) = \left\{ x \Big| f(x) \in f(E) \right\}$.

            Suppose that $E = f^{-1}(f(E))$ and $x,y \in X$.
            Let $x \in E$ and $y \notin E$, so $x \neq y$.
            Therefore, if $x \in f^{-1}(f(E))$ and $y \notin f^{-1}(f(E))$.
            This also means that $f(x) \in f(E)$.
            Additionally, $f(y) \notin f(E)$, because if $f(y) \in f(E)$, then by definition $y \in f^{-1}(f(E))$, which is a contradiction.
            Therefore, $x \neq y$ implies $f(x) \neq f(y)$.
            By contrapositive, this means that $f(x) = f(y)$ implies $x = y$.
            Ergo, $f$ is injective.
        \end{proof}

        Last, we consolidate.
        \begin{proof}
            We know that if $f$ is injective, then $E = f^{-1}(f(E))$.
            We also know that if $E = f^{-1}(f(E))$, then $f$ is injective.
            Therefore, $f$ is injective if and only if $E = f^{-1}(f(E))$.
        \end{proof}

    \subsection{Solution (d)}
        First we prove that $f(f^{-1}(E)) \subset E$.
        \begin{proof}
            Take $y \in f(f^{-1}(E))$.
            Therefore, there exists some element $x \in f^{-1}(E)$ such that $f(x) = y$.
            Since $x \in f^{-1}(E)$, then $f(x) \in E$.
            Since $y = f(x)$, $y \in E$ for any $y \in f(f^{-1}(E))$.
            Ergo, $f(f^{-1}(E)) \subset E$.
        \end{proof}

        Second, we prove that if $f$ is surjective, then $f(f^{-1}(E)) = E$.
        \begin{proof}
            Take any $x \in E$.
            Since $f$ is surjective, for every element $z \in E$, there is an element $y \in X$ such that $f(y) = z$.
            Therefore, there is and element $y \in X$ such that $f(y) = x$.
            By the definition of $f^{-1}(E)$, this means that $y \in f^{-1}(E)$.
            Furthermore, by the definition of $f(f^{-1}(E))$, since $y \in f^{-1}(E)$, $f(y) \in f(f^{-1}(E))$.
            Since $f(y) = x$, $x \in f(f^{-1}(E))$.
            Therefore, $x \in E$ implies $x \in f(f^{-1}(E))$.
            Ergo, $E \subset f(f^{-1}(E))$.

            We have established that $E \subset f(f^{-1}(E))$.
            We have established that $f(f^{-1}(E)) \subset E$.
            Ergo, $f(f^{-1}(E)) = E$.
        \end{proof}

        Third, we prove that if $f(f^{-1}(E)) = E$, then $f$ is surjective.
        \begin{proof}
            Suppose that for all $E$, $f(f^{-1}(E)) = E$.
            This means that if $y \in E$ then $y \in f(f^{-1}(E))$.
            Since $y \in f(f^{-1}(E))$, there exists $x \in f^{-1}(E)$ where $f(x) = y$.
            Since $f^{-1}(E) \subset X$, $x \in X$.
            Therefore, for any $y \in E$, there exists $x \in X$ such that $f(x) = y$.
            Ergo, $f$ is surjective within $E$ for all $E \subset Y$.

            Suppose that $E = Y$.
            That would mean that $f$ is surjective in $E$ and as such surjective in $Y$.
            Therefore, $f$ is surjective over $Y$.
        \end{proof}

        Last, we consolidate.
        \begin{proof}
            We know that if $f$ is surjective, then $E = f(f^{-1}(E))$.
            We also know that if $E = f(f^{-1}(E))$, then $f$ is surjective.
            Therefore, $f$ is surjective if and only if $E = f(f^{-1}(E))$.
        \end{proof}

\pagebreak
\section{Problem 4}
    Let $(X, d_X)$ and $(Y, d_Y)$ be metric spaces.
    \begin{enumerate}[label=(\alph*)]
        \item 
        A function $f : X \to Y$ (not necessarily continuous) is said to be \textit{locally bounded} if, for each $p \in X$, there is an $r > 0$ so that $f(B_r(p))$ is a bounded set in $Y$. Show that every locally bounded function on a \textit{compact} metric space is bounded (i.e. $f(X)$ is a bounded set in $Y$).

        \item 
        A function $f : X \to Y$ is said to be \textit{locally constant} if, for each $p \in X$, there is an $r > 0$ so that $f$ is constant on $B_r(p)$. Show that every locally constant function on a \textit{connected} metric space is constant. \textit{Hint:} show that $f^{-1}(\{y\})$ is clopen in $X$ for any $y \in Y$
        
    \end{enumerate}
    \subsection{Solution (a)}
        \begin{proof}
            Let $f: X \to Y$ be locally bounded and let $(X, d_X)$ be compact.
            Since $(X, d_X)$ is compact, every open cover of $X$ has a finite subcover.
            By the definition of a locally bounded function, for each $p \in X$, there is an $r_p > 0$ so that $f(B_{r_p}(p))$ is bounded in $Y$.
            Create a collection out of every $B_{r_p}(p)$ and call it $P$.
            Since every open ball contains its center and open balls are open, the collection is an open cover of $X$.
            Since $X$ is compact, there is a cover of $X$ made up of a finite number of elements in $P$, which we can call $F$.
            
            Create a new set called $G$ defined as $G = \left\{ f(S) \right\}_{S \in F}$.
            Since every element in $F$ is an open ball, applying $f$ to every open ball returns a bounded set in $Y$, and $F$ is finite, $G$ is a finite collection of bounded sets in $Y$.
            Since the union of a finite number of bounded sets is bounded\footnote{Prove or cite a source.}, the union of all elements in $G$, $H = \bigcup_{S \in G} S$, is a bounded set.
            Furthermore, since $H$ is a collection of bounded sets in $Y$, $H \subset Y$.

            Now, we prove that if $f(x) \in f(X)$, then $f(x) \in H$.
            Take any element $f(x) \in f(X)$.
            By the definition of $f(X)$, $x \in X$.
            Since $F$ is a cover of $X$, there exists $T \in F$ such that $x \in T$.
            Therefore, $f(x) \in f(T)$.
            Since $H = \bigcup_{S \in G} S = \bigcup_{S \in F} f(S)$, $f(T) \subset H$.
            $f(x) \in f(T)$ and $f(T) \subset H$ means $f(x) \in H$.
            Therefore, $f(x) \in f(X)$ implies $f(x) \in H$.
            Therefore, $f(X) \subset H$.

            Lastly, a few facts to put everything together.
            Since $H$ is bounded and $f(X) \subset H$, $f(X)$ is bounded.
            Since $H \subset Y$ and $f(X) \subset H$, $f(X) \subset Y$.
            Therefore, $f(X)$ is a bounded set in $Y$.
            Ergo, every locally bounded function on a compact metric space is bounded.
        \end{proof}

    \subsection{Solution (b)}
        \begin{proof}[Proof by Contradiction]
            Let $f: X \to Y$ be locally constant and let $(X, d_X)$ and $(Y, d_Y)$ be metric spaces.
            Furthermore, let $(X, d_X)$ be connected, meaning the only clopen sets in $(X, d_X)$ are $\varnothing$ and $X$.
            
            Suppose for contradiction that $f$ is not constant. 
            This means that there exists a set $V \subset Y$ with more than one element where $V = f(X)$.
            For all $v \in V$, let $X_v$ be the set of elements $p \in X$ where $f(p) = v$.
            Since $f$ is locally constant, for each $p \in X$, there is an $r > 0$ such that for all $q \in B_r(p)$, $f(p) = f(q)$.
            If $p \in X_v$, then $q \in X_v$ and by extension $B_r(p) \subset X_v$.
            Therefore, for all $v \in V$, $X_v$ is open.

            Since $f$ is a well-defined function, for every $p \in X$, there exists $f(p) \in Y$.
            Therefore, for every $p \in X$, there exists $v \in V$ where $f(p) = v$.
            Therefore, for every $p \in X$, there exists $v \in V$ such that $p \in X_v$.
            Therefore, $X = \bigcup_{v \in V} X_v$ and for any $w \in V$, $X_w^c = X \backslash X_w = \bigcup_{v \in V \backslash \{w\}} X_v$.

            Take any $w \in V$.
            Since $V = f(X)$, $w \in f(X)$, meaning there is some element $x \in X$ where $f(x) = w$.
            Therefore, $X_w \neq \varnothing$.
            Furthermore, $f$ is not constant, so $X_w \neq X$.
            We already proved that for every $v \in V$, $X_v$ is open.
            Therefore, since the union of open sets is open, $\bigcup_{v \in V \backslash \{w\}} X_v$ is also open.
            This means that $X_w^c$ is open.
            Since a set is closed if and only if its complement is open, $X_w$ is closed.
            Therefore, $X_w$ is both open and closed, or clopen.
            This contradicts the fact that $X$ is connected, since the only clopen subsets of connected set $X$ are $\varnothing$ and $X$.
            Ergo, $f(X)$ is constant.
        \end{proof}

    \pagebreak
\section{Problem 5}
    Using the $\epsilon-\delta$ definition of continuity directly, prove that the following functions are continuous.
    \begin{enumerate}[label=(\alph*)]
        \item 
        $f : \mathbb{R} \to \mathbb{R}$; \quad $f(x) = 104x$.

        \item 
        $f : \mathbb{R} \to \mathbb{R}$; \quad $f(x) = x^2$.
        
    \end{enumerate}

    \subsection{Solution (a)}
        \begin{proof}
            For a function to be continuous, it must be continuous at every point in its domain.
            Recall the $\epsilon-\delta$ definition of continuous at a point: for all $\varepsilon > 0$, there exists some $\delta > 0$ such that $d_X(x,p) < \delta$ implies $d_Y(f(x),f(p)) < \varepsilon$.
            In the real numbers and the usual metric, this means that $\abs{x - p} < \delta$ implies that $\abs{f(x) - f(p)} < \varepsilon$.

            Let $\varepsilon > 0$ and $\delta = \frac{\varepsilon}{104} > 0$.
            Suppose that $\abs{x - p} < \delta$.
            This means that $\abs{x - p} < \frac{\varepsilon}{104}$.
            Multiplying everything by $104$, we find that $\varepsilon > 104 \abs{x - p} = \abs{104x - 104p}$.
            Since $f(x) = 104x$, $\abs{f(x) - f(p)} < \varepsilon$.
            Putting things together, this means that for any $\varepsilon > 0$, there exists some $\delta > 0$ such that $\abs{x - p} < \delta$ implies $\abs{f(x) - f(p)} < \varepsilon$.
            Ergo, $f(x) = 104x$ is continuous.
        \end{proof}

    \subsection{Solution (b)}
        \begin{proof}
            Let $\varepsilon > 0$.
            Take any $x, p \in \mathbb{R}$.
            By the definition of the absolute value, $\abs{x + p} \geq 0$.
            There are two conditions from here, that either $x = p$ or $x \neq p$.

            Suppose $x = p$.
            This means that $\abs{x - p} = 0$, so for any $\delta > 0$, $\delta > \abs{x - p}$.
            Furthermore, since $\abs{f(x) - f(p)} = \abs{x^2 - p^2} = \abs{x - p} \abs{x + p}$, $\abs{x^2 - p^2} = \abs{x - p} \abs{x + p} = 0 \abs{x + p} = 0$.
            Therefore, $\varepsilon > \abs{f(x) - f(p)}$, so if $x = p$, the $\epsilon-\delta$ definition is satisfied.

            Suppose $x \neq p$, so $\abs{x + p} \neq 0$, so $\abs{x + p} > 0$ by the definition of the absolute value.
            This means $\abs{x + p} > 0$.
            Let $\delta = \frac{\varepsilon}{\abs{x + p}}$ and suppose that $\abs{x - p} < \delta$.
            Therefore, $\abs{x - p} < \frac{\varepsilon}{\abs{x + p}}$
            Multiply both sides by $\abs{x + p}$ to find that $\abs{x - p} \abs{x + p} < \varepsilon$.
            Therefore, $\abs{(x - p)(x + p)} < \varepsilon$.
            Therefore, $\abs{x^2 - p^2} < \varepsilon$.
            This means $\abs{f(x) - f(p)} < \varepsilon$.
            Therefore, the $\epsilon-\delta$ definition is satisfied.

            The $\epsilon-\delta$ definition is satisfied in all cases.
            Ergo, $f(x) = x^2$ is continuous.
        \end{proof}

    \pagebreak
\section{Problem 6}
    Let $f : \mathbb{R} \to \mathbb{R}$ be defined by $f(x) = \begin{cases} x, & \text{ if } x \in \mathbb{Q}\\ 0, &\text{ if } x \notin \mathbb{Q}\end{cases}.$ Using the $\epsilon-\delta$ definition of continuity,
    \begin{enumerate}[label=(\alph*)]
        \item 
        prove that $f$ is not continuous at any point $x \neq 0$;
        
        \item
        prove that $f$ \textit{is} continuous at $x = 0$.
        
    \end{enumerate}

    \subsection{Solution (a)}
        This is an altered version of Thomae's Function.
        We consider this in two cases: $x \in \mathbb{Q} \backslash \{0\}$ and $x \neq \mathbb{Q}$.

        First we prove that $f$ is not continuous at $x$ if $x \in \mathbb{Q} \backslash \{0\}$.
        \begin{proof}
            Take $\varepsilon = \frac{\abs{f(x)}}{2}$.
            Take any $\delta > 0$ such that $\abs{x - p} < \delta$.
            Since the irrational numbers are dense in the real numbers, between any two real numbers, there is an irrational number.
            This means that there is an irrational number $p$ where $x - \delta < p < x + \delta$.
            This means that $-\delta < p - x < \delta$, and by extension $\abs{p - x} = \abs{x - p} < \delta$.

            Now take $\abs{f(x) - f(p)}$.
            Since $x \in \mathbb{Q}$ and $p \notin \mathbb{Q}$.
            Therefore, $\abs{f(x) - f(p)} = \abs{x - 0} = \abs{x}$.
            Therefore, since $\varepsilon = \frac{\abs{f(x)}}{2} = \frac{\abs{x}}{2}$, $\abs{x} = \abs{f(x) - f(p)} \geq \varepsilon$.
            Ergo, by the $\epsilon-\delta$ definition, $f$ is not continuous at $x \in \mathbb{Q} \backslash \{0\}$.
        \end{proof}

        Next, we prove that $f$ is not continuous at $x$ if $x \in \mathbb{Q}$.
        \begin{proof}
            Take any $\delta > 0$ such that $\abs{x - p} < \delta$.
            Since the rational numbers are dense in the real numbers, there is a non-zero rational number $p$ where $x - \delta < p < x + \delta$.
            Therefore, $\abs{x - p} < \delta$.

            Now, look at the epsilon-side.
            Since $p$ is rational but $x$ is irrational, $\abs{f(x) - f(p)} = \abs{p - 0} = \abs{p}$.
            If $\varepsilon > \abs{p} > 0$, then $\varepsilon > \abs{f(x) - f(p)}$.
            However, if $0 < \varepsilon \leq \abs{p}$, then $\varepsilon \leq \abs{f(x) - f(p)}$.
            Ergo, by the $\epsilon-\delta$ definition of continuous, $f$ is not continuous at $x$ if $x \in \mathbb{Q}$.
        \end{proof}

        If necessary, we consolidate this.
        \begin{proof}
            We know that $f$ is not continuous at $x$ if $x \in \mathbb{Q}$.
            We know $f$ is not continuous at $x$ if $x \in \mathbb{Q} \backslash \{0\}$.
            Therefore, $f$ is not continuous at $x$ if $x \in \mathbb{R} \backslash \{0\}$.
            Ergo, $f$ is not continuous at any point $x \neq 0$.
        \end{proof}

    \subsection{Solution (b)}
        \begin{proof}
            Let $\varepsilon > 0$.
            Suppose that $\delta = \varepsilon$.
            Assume that $\abs{x - p} < \delta$.
            This means that $\abs{p - x} < \delta$.
            Since $x = 0$, $\abs{p} < \delta$.
            Since  $\delta = \varepsilon$, this also means $\abs{p} < \varepsilon$.

            Now, we apply this to the function.
            Suppose that $p \notin \mathbb{Q}$.
            This means that $\abs{x - p} = \abs{0 - 0} = \abs{0} = 0$.
            By definition, we know that this works, since $\varepsilon > 0$.
            Suppose that $p \in \mathbb{Q}$.
            This means that $\abs{x - p} = \abs{p - 0} = \abs{p}$.
            Earlier in the proof, we established that $\abs{p} < \varepsilon$, so it holds.

            Ergo, $f$ is continuous at $x = 0$.
        \end{proof}
\end{document}